% Part: applied-modal-logics
% Chapter: epistemic-logic
% Section: bisimulations

\documentclass[../../../include/open-logic-section]{subfiles}

\begin{document}

\olfileid{aml}{el}{bsd}

\olsection{في المحاكاة الثنائية}

ويبقى في قوة اللغة المعرفية سؤال عن صلة النموذج بالصيغ الصادقة فيه. فقد
دلت نتائج مطابقة الأطر على أن صدق بعض الصيغ صدقًا عامًا في إطار يستلزم
اتصافه بخواص معينة. ولكن هل تميز لغتنا الجهوية بين عالم ذي سهم انعكاسي
وبين سلسلة لا تنتهي من العوالم، ينفذ كل واحد منها إلى الذي يليه؟ وبعبارة
أخرى: هل توجد صيغة~${\OLId{latin-looped}{A}}$ تصدق في أحد الموضعين دون الآخر؟

وللجواب نعرّف بين النماذج العلائقية علاقة تسمى المحاكاة الثنائية؛ وهي
تضبط اشتراكها في البنية التي تستطيع اللغة المعرفية التعبير عنها، ومن ثم
تعطي معنى للتكافؤ بينها.

\begin{defn}[المحاكاة الثنائية]
ليكن ${\OLId{latin-looped}{M}}_{\OLMathNumeral{1}} = \tuple{{\OLId{latin-looped}{W}}_{\OLMathNumeral{1}},({\OLId{latin-looped}{R}}_{{\OLMathNumeral{1}},{\OLId{latin-ordinary}{a}}})_{{\OLId{latin-ordinary}{a}} \in {\OLId{latin-looped}{G}}},{\OLId{latin-looped}{V}}_{\OLMathNumeral{1}}}$ و
${\OLId{latin-looped}{M}}_{\OLMathNumeral{2}} = \tuple{{\OLId{latin-looped}{W}}_{\OLMathNumeral{2}},({\OLId{latin-looped}{R}}_{{\OLMathNumeral{2}},{\OLId{latin-ordinary}{a}}})_{{\OLId{latin-ordinary}{a}} \in {\OLId{latin-looped}{G}}},{\OLId{latin-looped}{V}}_{\OLMathNumeral{2}}}$ نموذجين علاقيين، ولتكن
$\mathcal{R} \subseteq {\OLId{latin-looped}{W}}_{\OLMathNumeral{1}} \times {\OLId{latin-looped}{W}}_{\OLMathNumeral{2}}$ علاقة ثنائية. نقول إن
$\mathcal{R}$ \emph{محاكاة ثنائية} إذا تحققت الشروط الآتية لكل
$\tuple{{\OLId{latin-ordinary}{w}}_{\OLMathNumeral{1}}, {\OLId{latin-ordinary}{w}}_{\OLMathNumeral{2}}} \in \mathcal{R}$:

\begin{enumerate}
  \item ${\OLId{latin-ordinary}{w}}_{\OLMathNumeral{1}} \in {\OLId{latin-looped}{V}}_{\OLMathNumeral{1}}({\OLId{latin-ordinary}{p}})$ إذا وفقط إذا كان ${\OLId{latin-ordinary}{w}}_{\OLMathNumeral{2}} \in {\OLId{latin-looped}{V}}_{\OLMathNumeral{2}}({\OLId{latin-ordinary}{p}})$ لكل
  !!{propositional variable}s~${\OLId{latin-ordinary}{p}}$.
  \item لكل فاعل ${\OLId{latin-ordinary}{a}} \in {\OLId{latin-looped}{G}}$ ولكل عالم ${\OLId{latin-ordinary}{v}}_{\OLMathNumeral{1}} \in {\OLId{latin-looped}{W}}_{\OLMathNumeral{1}}$، إذا كان
  ${\OLId{latin-looped}{R}}_{{\OLMathNumeral{1}},{\OLId{latin-ordinary}{a}}} {\OLId{latin-ordinary}{w}}_{\OLMathNumeral{1}} {\OLId{latin-ordinary}{v}}_{\OLMathNumeral{1}}$، وُجد ${\OLId{latin-ordinary}{v}}_{\OLMathNumeral{2}} \in {\OLId{latin-looped}{W}}_{\OLMathNumeral{2}}$ بحيث
  ${\OLId{latin-looped}{R}}_{{\OLMathNumeral{2}},{\OLId{latin-ordinary}{a}}} {\OLId{latin-ordinary}{w}}_{\OLMathNumeral{2}} {\OLId{latin-ordinary}{v}}_{\OLMathNumeral{2}}$ و$\tuple{{\OLId{latin-ordinary}{v}}_{\OLMathNumeral{1}}, {\OLId{latin-ordinary}{v}}_{\OLMathNumeral{2}}} \in
  \mathcal{R}$.
  \item لكل فاعل ${\OLId{latin-ordinary}{a}} \in {\OLId{latin-looped}{G}}$ ولكل عالم ${\OLId{latin-ordinary}{v}}_{\OLMathNumeral{2}} \in {\OLId{latin-looped}{W}}_{\OLMathNumeral{2}}$، إذا كان
   ${\OLId{latin-looped}{R}}_{{\OLMathNumeral{2}},{\OLId{latin-ordinary}{a}}} {\OLId{latin-ordinary}{w}}_{\OLMathNumeral{2}} {\OLId{latin-ordinary}{v}}_{\OLMathNumeral{2}}$، وُجد ${\OLId{latin-ordinary}{v}}_{\OLMathNumeral{1}} \in {\OLId{latin-looped}{W}}_{\OLMathNumeral{1}}$ بحيث
   ${\OLId{latin-looped}{R}}_{{\OLMathNumeral{1}},{\OLId{latin-ordinary}{a}}} {\OLId{latin-ordinary}{w}}_{\OLMathNumeral{1}} {\OLId{latin-ordinary}{v}}_{\OLMathNumeral{1}}$ و$\tuple{{\OLId{latin-ordinary}{v}}_{\OLMathNumeral{1}}, {\OLId{latin-ordinary}{v}}_{\OLMathNumeral{2}}} \in
   \mathcal{R}$.
\end{enumerate}

إذا وُجدت محاكاة ثنائية بين ${\OLId{latin-looped}{M}}_{\OLMathNumeral{1}}$ و${\OLId{latin-looped}{M}}_{\OLMathNumeral{2}}$ تربط العالمين ${\OLId{latin-ordinary}{w}}_{\OLMathNumeral{1}}$
و${\OLId{latin-ordinary}{w}}_{\OLMathNumeral{2}}$، جاز لنا أيضًا أن نكتب
$\tuple{{\OLId{latin-looped}{M}}_{\OLMathNumeral{1}}, {\OLId{latin-ordinary}{w}}_{\OLMathNumeral{1}}} \leftrightarroweq \tuple{{\OLId{latin-looped}{M}}_{\OLMathNumeral{2}}, {\OLId{latin-ordinary}{w}}_{\OLMathNumeral{2}}}$، وأن نسمي
$\tuple{{\OLId{latin-looped}{M}}_{\OLMathNumeral{1}}, {\OLId{latin-ordinary}{w}}_{\OLMathNumeral{1}}}$ و$\tuple{{\OLId{latin-looped}{M}}_{\OLMathNumeral{2}}, {\OLId{latin-ordinary}{w}}_{\OLMathNumeral{2}}}$ \emph{متكافئين بالمحاكاة
الثنائية}.
\end{defn}
% تصحيح دلالي (0488-I1): صُرّح بعائلة علاقات النفاذ المفهرسة لكل فاعل في كلا النموذجين.

ولكل بند من بنود المحاكاة الثنائية عمله: فالبند~\OLClassicalDigits{1} يجعل العالمين المتكافئين
يرضيان الصيغ نفسها الخالية من المؤثرات الجهوية، لأنه يوجب اتفاقهما في جميع
ال!!{propositional variable}s. وأما البندان الآخران، ويسميان على الترتيب
«الذهاب» و«الإياب»، فيحفظان بنية علاقات الإتاحة.

\begin{thm}
إذا كان $\tuple{{\OLId{latin-looped}{M}}_{\OLMathNumeral{1}}, {\OLId{latin-ordinary}{w}}_{\OLMathNumeral{1}}} \leftrightarroweq \tuple{{\OLId{latin-looped}{M}}_{\OLMathNumeral{2}}, {\OLId{latin-ordinary}{w}}_{\OLMathNumeral{2}}}$، فلكل
!!{formula}~$!A$ يكون $\mSat{{\OLId{latin-looped}{M}}_{\OLMathNumeral{1}}}{!A}[{\OLId{latin-ordinary}{w}}_{\OLMathNumeral{1}}]$ إذا وفقط إذا كان
$\mSat{{\OLId{latin-looped}{M}}_{\OLMathNumeral{2}}}{!A}[{\OLId{latin-ordinary}{w}}_{\OLMathNumeral{2}}]$.
\end{thm}

\begin{figure}
  \begin{center}
    \begin{tikzpicture}[modal]
      \node[world] (w1) {${\OLId{latin-ordinary}{w}}_{\OLMathNumeral{1}}$}; 
      \node[world] (w2) [above left=of w1]{${\OLId{latin-ordinary}{w}}_{\OLMathNumeral{2}}$}; 
      \node[world] (w3) [above right=of w1] {${\OLId{latin-ordinary}{w}}_{\OLMathNumeral{3}}$};
      \draw[<->] (w1) to node {${\OLId{latin-ordinary}{a}}$} (w2);
      \draw[->,loop below] (w1) to node {${\OLId{latin-ordinary}{a}}$} (w1);
      \draw[<->] (w1) to [swap] node {${\OLId{latin-ordinary}{a}}$} (w3);
      \draw[->,loop above] (w2) to node {${\OLId{latin-ordinary}{a}}$} (w2) ;
      \draw[->,loop above] (w3) to node {${\OLId{latin-ordinary}{a}}$} (w3) ;
      
      \node[world](v1)[right of=w1, xshift=1.5in]{${\OLId{latin-ordinary}{v}}_{\OLMathNumeral{1}}$};
      \node[world](v2)[above of=v1]{${\OLId{latin-ordinary}{v}}_{\OLMathNumeral{2}}$};
      \draw[<->] (v1) to node {${\OLId{latin-ordinary}{a}}$} (v2);
      \draw[->,loop below] (v1) to node {${\OLId{latin-ordinary}{a}}$} (v1);
      \draw[->,loop above] (v2) to node {${\OLId{latin-ordinary}{a}}$} (v2) ;

      \draw[-,dotted] (w1) to (v1) ;
      \draw[-,dotted,bend left=45] (w2) to (v2) ;
      \draw[-,dotted,bend left=30] (w3) to (v2) ;
    \end{tikzpicture}
  \end{center}
  \caption{نموذجان متكافئان بالمحاكاة الثنائية.}
  \ollabel{fig:bisimilar}
\end{figure}

ليس النموذجان المرسومان في \olref{fig:bisimilar} متماثلين، ومع ذلك تصل
محاكاة ثنائية العالم~${\OLId{latin-ordinary}{w}}_{\OLMathNumeral{1}}$ بالعالم~${\OLId{latin-ordinary}{v}}_{\OLMathNumeral{1}}$. وهي تصل كذلك كلًا من~${\OLId{latin-ordinary}{w}}_{\OLMathNumeral{2}}$
و${\OLId{latin-ordinary}{w}}_{\OLMathNumeral{3}}$ بــ${\OLId{latin-ordinary}{v}}_{\OLMathNumeral{2}}$؛ إذ لا تستطيع لغتنا الجهوية أن تعبّر عن فرق يميز أحدهما
من الآخر. ولو أرضى~${\OLId{latin-ordinary}{w}}_{\OLMathNumeral{2}}$ و${\OLId{latin-ordinary}{w}}_{\OLMathNumeral{3}}$ !!{propositional variable}s مختلفة
لتغير الحكم.

\end{document}
